% SPDX-License-Identifier: Apache-2.0
\section{A Debugger From the Cable}
\label{sec:x-dtm}

The RISC-V debug specification puts a debugger's access to a hart
behind a small space of registers, the debug module interface, which
a transport reaches from a cable. On an FPGA the JTAG port belongs to
the device, so a design reaches it through \code{BSCANE2}, one of the
user scan chains of the device's own TAP. \code{Dtm} is the
specification's transport on that chain (issue 154). It is what lets
a stock OpenOCD, and gdb through it, reach the debug module the board
already has, with nothing written for either of them.

OpenOCD communicates with it through its BSCAN tunnel. It sets the device's
instruction register to \code{USER4} and encapsulates every scan the
transport's own TAP would have processed inside one data scan of that
chain: a bit that says instruction or data, seven bits of width, the
payload, and three clocks after it. The payload is one bit longer than
the register, because what comes back is late by a clock. The
transport has the specification's registers: its \code{IDCODE},
\code{dtmcs}, and \code{dmi}, the 41 bits of an access, which go out on
\code{req} and come back on \code{ans}. An access still in flight when
the next scan comes says busy, which sticks until \code{dtmcs} is
written to clear it, as the specification has it. All of it runs on
\code{Tck}, the cable's clock, which is a clock of its own in this
language like any other.

The example plays the scans OpenOCD makes, field for field from its
source, against a model of the debug module that answers each access
two edges after it comes: the \code{IDCODE}, \code{dtmcs}, then a
write of a register and a read of it back. The driver changes the
pins on the falling edge, as a JTAG host does, so that no edge sees
them change under it. The netlist is checked against the run under nvc
and Verilator.

\begin{figure}[htbp]
\centering
\input{dtm_timing.tex}
\caption{One tunneled scan of the instruction register, which selects
\code{dtmcs}: a capture, \code{phase} \code{01} for the seven bits of
width, counted in \code{cnt}, \code{10} for the five bits of the
register, \code{11} for the clocks after it, and on the update
\code{ir} goes from \code{1}, \code{IDCODE}, to \code{10},
\code{dtmcs}. The update at the left is the scan before it.}
\label{fig:x-dtm}
\end{figure}

\lstinputlisting[caption={\code{ex\_dtm.rs}. Run with \code{bazel run //lib/examples:ex\_dtm}.},label={lst:x-dtm}]{lib/examples/ex_dtm.rs}

\lstinputlisting[style=ascii,caption={What \code{ex\_dtm} printed when the build ran it: the transport's \code{IDCODE}, \code{dtmcs}, and the word read back through \code{dmi}, then the netlist.},label={lst:x-dtm-out}]{out_dtm.txt}
